\chapter{El código}
\label{ap:codigo}

Todo el material numérico de este manual se produjo con una implementación
propia, escrita desde cero y sin dependencias más allá de \texttt{numpy},
\texttt{scipy} y \texttt{matplotlib}. La prioridad fue la legibilidad: el código
está pensado para leerse junto con los capítulos \ref{cap:bloque} a
\ref{cap:dag}.

\section{Estructura}

\begin{center}\small
\begin{tabular}{@{}lp{0.62\textwidth}@{}}
\toprule
\texttt{ssj/grids.py} & grillas y discretización de Rouwenhorst\\
\texttt{ssj/interpolate.py} & interpolación, lotería de Young, pasos hacia
adelante y de expectativa\\
\texttt{ssj/het.py} & el bloque heterogéneo y el algoritmo fake news\\
\texttt{ssj/simple.py} & bloques simples, grafo acíclico, solución de equilibrio\\
\texttt{models/krusell\_smith.py} & el modelo de referencia\\
\texttt{models/banks.py} & el bloque bancario de la parte V\\
\texttt{models/banks\_ge.py} & su equilibrio general\\
\texttt{tests/test\_fakenews.py} & verificación contra el método directo\\
\texttt{figures/} & generación de todas las figuras\\
\bottomrule
\end{tabular}
\end{center}

\section{Cómo correrlo}

\begin{lstlisting}[language=bash,numbers=none,basicstyle=\ttfamily\footnotesize]
pip install numpy scipy matplotlib
cd code
python3 tests/test_fakenews.py     # verificacion (deberia decir TODO OK)
python3 figures/fig_metodo.py      # figuras 1 a 8
python3 figures/fig_bancos.py      # figuras 9 a 13
\end{lstlisting}

La verificación tarda unos segundos; las figuras del modelo bancario tardan unos
minutos la primera vez, porque resuelven el punto fijo del estado estacionario, y
después usan un caché en disco.

\section{El núcleo del algoritmo}

Vale la pena leer estas treinta líneas con la receta del capítulo
\ref{cap:fakenews} al lado. Son la traducción literal de los cuatro pasos.

\paragraph{Paso 1: una sola iteración hacia atrás}

\begin{lstlisting}
def _backward_fakenews(self, ss, shocked, T, h):
    """Una sola iteracion hacia atras entrega TODAS las columnas."""
    base = {k: ss[k] for k in self.inputs}
    curlyY = {o: np.empty((2, T)) for o in self.outputs}
    curlyD = np.empty((2, T, self.ne, self.na))

    for sgn_idx, sgn in enumerate((+1.0, -1.0)):   # diferenciacion de dos lados
        V_next = ss["V"]
        for s in range(T):
            args = dict(base)
            if s == 0:
                args[shocked] = base[shocked] + sgn * h
            V_next, pol, out = self._step(V_next, **args)
            for o in self.outputs:
                curlyY[o][sgn_idx, s] = np.vdot(out[o], ss["D"])   # dY_0^s
            i, p = lottery(pol, self.grid)
            curlyD[sgn_idx, s] = self._forward(ss["D"], i, p)      # dD_1^s

    dY = {o: (curlyY[o][0] - curlyY[o][1]) / (2 * h) for o in self.outputs}
    dD = (curlyD[0] - curlyD[1]) / (2 * h)
    return dY, dD
\end{lstlisting}

El punto que hay que ver: el bucle sobre \texttt{s} no reinicia nada. La
variable \texttt{V\_next} se arrastra de una iteración a la siguiente, de modo
que en la iteración $s$ la función de continuación ya incorpora el efecto de un
choque que ocurre $s$ períodos en el futuro. Eso es el lema
\ref{lema:politicas} hecho código.

\paragraph{Pasos 2 a 4: expectativas, matriz $\F$ y acumulación}

\begin{lstlisting}
def _expectations(self, ss, T):
    """E_t = Lambda^t y_ss: T-1 aplicaciones de una matriz fija."""
    curlyE = {}
    for o in self.outputs:
        E = np.empty((T - 1, self.ne, self.na))
        E[0] = ss["outcomes"][o]
        for t in range(1, T - 1):
            E[t] = self._expect(E[t - 1], ss["i"], ss["p"])
        curlyE[o] = E
    return curlyE

def jacobians(self, ss, shocked_inputs=None, T=300, h=1e-4, return_F=False):
    shocked_inputs = shocked_inputs or self.inputs
    curlyE = self._expectations(ss, T)                       # paso 2
    dY, dD = {}, {}
    for i in shocked_inputs:
        dY[i], dD[i] = self._backward_fakenews(ss, i, T, h)  # paso 1

    Fs, Js = {}, {}
    for o in self.outputs:
        Emat = curlyE[o].reshape(T - 1, -1)
        for i in shocked_inputs:
            F = np.empty((T, T))                             # paso 3
            F[0, :] = dY[i][o]
            F[1:, :] = Emat @ dD[i].reshape(T, -1).T
            Fs[(o, i)] = F
            J = F.copy()                                     # paso 4
            for t in range(1, T):
                J[t, 1:] += J[t - 1, :-1]
            Js[(o, i)] = J
    return (Js, Fs) if return_F else Js
\end{lstlisting}

El paso 3 es una sola multiplicación de matrices de $(T-1)\times n_g$ por
$n_g\times T$, y el paso 4 es una recursión de $T$ líneas sobre una matriz
$T\times T$. Ahí está todo el algoritmo.

\section{Los dos operadores que hay que no confundir}

\begin{lstlisting}
def forward_step(D, Pi, i, p):
    """D_{t+1} = Lambda' D_t: primero la loteria, luego Markov exogeno."""
    return Pi.T @ forward_policy(D, i, p)

def expectation_step(E, Pi, i, p):
    """E_t = Lambda E_{t-1}: primero Markov exogeno, luego la loteria.

    Notese el orden INVERSO. Es la transpuesta del mismo operador: si
    Lambda' = Pi' L'  entonces  Lambda = L Pi. Confundir este orden produce
    jacobianos casi correctos, y diagnosticarlo cuesta horas.
    """
    Etilde = Pi @ E
    out = np.empty_like(E)
    for e in range(E.shape[0]):
        out[e] = p[e] * Etilde[e, i[e]] + (1 - p[e]) * Etilde[e, i[e] + 1]
    return out
\end{lstlisting}

\section{La verificación}

\begin{lstlisting}
def test_budget_and_signs():
    het = make_household(n_a=100, amax=150)
    ss = household_ss(het, r=0.01, w=1.0, eis=0.5, target_A=3.0)
    J = het.jacobians(ss, ["r", "w"], T=30)

    # restriccion presupuestaria agregada en la fecha 0:
    #   dC_0 + dA_0 = dw_0 * E[e]
    lhs = J[("C", "w")][0, 0] + J[("A", "w")][0, 0]
    assert abs(lhs - 1.0) < 1e-5

    # una NOTICIA sobre el futuro no cambia los recursos de hoy:
    #   dC_0 + dA_0 = 0  para s >= 1
    for s in (1, 5, 10):
        assert abs(J[("C", "w")][0, s] + J[("A", "w")][0, s]) < 1e-6
\end{lstlisting}

Esta prueba no mira ninguna magnitud económica: solo verifica contabilidad. Por
eso es tan útil. Un error de alineación temporal, de orden de operadores o de
índices rompe la contabilidad mucho antes de volver implausible el resultado
económico.

\section{Resultados de la verificación}

Al correr \texttt{tests/test\_fakenews.py} sobre el bloque de hogares de
Krusell--Smith con $T=40$ y 100 puntos de activos:

\begin{center}\small
\begin{tabular}{@{}lrr@{}}
\toprule
& \textbf{error absoluto} & \textbf{error relativo}\\
\midrule
$\J^{A,r}$ & $3.3\times10^{-8}$ & $5.5\times10^{-9}$\\
$\J^{A,w}$ & $1.5\times10^{-8}$ & $1.7\times10^{-8}$\\
$\J^{C,r}$ & $1.2\times10^{-9}$ & $3.6\times10^{-9}$\\
$\J^{C,w}$ & $1.4\times10^{-8}$ & $9.7\times10^{-8}$\\
\bottomrule
\end{tabular}
\end{center}

Y sobre el bloque bancario de la parte V, con 78 pares insumo--producto, el peor
error relativo es $2.9\times10^{-8}$.
